\documentclass{report}
\usepackage{amsmath,amssymb}
\setlength{\parindent}{0mm}
\setlength{\parskip}{1em}
\begin{document}
\begin{center}
\rule{6in}{1pt} \
{\large Ty Thompson \\
{\bf Efficient Uncertainty Quantification for a Stochastic Ginzburg-Landau Model}}

Colorado School of Mines \\ Golden CO
\\
{\tt tythomps@mines.edu}\\
Mahadevan Ganesh\end{center}

We focus on the efficient simulation of nondeterministic critical
phenomena in the Ginzburg-Landau (GL) model for superconductivity.
Deterministic GL is widely used to study the formation of vortex
configurations
in thin superconductors. When such phenomena are essentially
nondeterministic, the Langevin version of the time dependent GL problem
applies, and simulation presents a significant computational challenge.

To investigate nondeterministic dynamics, we work on 2-manifolds having
rotational symmetry about the axis of a constant magnetic field, and
consider an ideal (superconductor or normal) initial state. Using highly
efficient deterministic schemes, we first motivate the stochastic
extension, and demonstrate an ad hoc approach for simulating the
evolution of dense, locally stable vortex configurations. For stochastic
simulations,
we identified an opportunity to use a spectral Galerkin approach.
Building upon a linearized Crank-Nicolson scheme, we used generalized
spectral decompositions to reliably achieve a reduction in
dimensionality. The
efficiency of the method is examined through comparisons with monte carlo
estimators, and our efforts at qualifying the perturbation approach are
discussed.


\end{document}
